% Notas de aula — Capítulo 11: Circulação Geral
% Curso: Meteorologia Prática (baseado em R. Stull, Practical Meteorology, CC BY-NC-SA 4.0)
% Caixas verdes = conteúdo literal do quadro; linhas verdes = encenação.
\documentclass[11pt]{article}
\usepackage{../comum/preambulo}
\graphicspath{{../figuras/cap11/}}

\title{Capítulo 11 --- Circulação Geral\\
{\large Notas de aula (roteiro de quadro-negro)}}
\author{Curso de Meteorologia Prática}
\date{}

\begin{document}
\maketitle
\thispagestyle{fancy}

\noindent\textbf{Plano sugerido:} 2 aulas de 2\,h.
\emph{Aula 1:} \S\ref{sec:motor}--\S\ref{sec:hadley} (o desequilíbrio
radiativo, o transporte exigido, célula de Hadley e jato subtropical).
\emph{Aula 2:} \S\ref{sec:termico}--\S\ref{sec:rossby} (vento térmico,
jato polar, ondas de Rossby, os cinturões de pressão).

\section{O motor: aquecimento diferencial}\label{sec:motor}

\begin{noquadro}[abertura]
\textbf{Cap.~11 --- Circulação Geral}\\[2pt]
trópicos: absorvem $>$ emitem \qquad polos: emitem $>$ absorvem\\
nenhuma latitude esquenta para sempre $\Rightarrow$ atmosfera $+$
oceanos \textbf{transportam} o excesso\\[2pt]
a área entre as duas curvas ($Q_{abs}$ e $OLR$) \emph{é} a meteorologia.
\end{noquadro}

\quadro{Desenhar os dois perfis por latitude: insolação absorvida (pico
no equador \javimos{2}{insolação}) e OLR quase plana. Perguntar: por que
a OLR é tão plana? (Porque a circulação já espalhou a energia — o
gráfico observado é o \emph{resultado} do transporte.)}

Este capítulo cobra as promessas de meia dúzia de capítulos: o
desequilíbrio veio do Cap.~2 \javimos{2}{balanço radiativo}; as
ferramentas de vento, do Cap.~10 \javimos{10}{geostrofia e gradiente};
a espessura proporcional à temperatura, do Cap.~1
\javimos{1}{hipsométrica} — e ela será a chave do vento térmico.

\begin{formadif}[o transporte exigido]
Em regime estacionário, o transporte meridional através da latitude
$\phi$ é a integral do desequilíbrio ao sul (HS) daquela latitude:
\[ T(\phi) = \int_{-90°}^{\phi} \bigl[Q_{abs}(\phi') - OLR(\phi')\bigr]\,
   2\pi R^2\cos\phi'\,d\phi'. \]
O teste de consistência é $T(90°) = 0$ — o planeta fecha o balanço
(T1).
\end{formadif}

\begin{noquadro}[o número para levar]
pico de $\sim$5--6\,PW perto de $35°$\\
($\sim$4\,PW na atmosfera, o resto no oceano)\\[2pt]
1\,PW $= 10^{15}$\,W $\simeq$ 70 vezes o consumo elétrico da humanidade
\end{noquadro}

Cálculo completo no Caso~1 do notebook; o N1 mostra que pôr gelo
(albedo alto) nos polos \emph{aumenta} o transporte exigido — um gostinho
da retroalimentação do gelo \veremos{21}{feedback de albedo}.

\section{Célula de Hadley e o jato subtropical}\label{sec:hadley}

\quadro{Desenhar a célula em corte meridional: sobe na ZCIT, escoa em
altitude, desce nos subtrópicos ($\sim30°$), retorna como alísios.
Marcar as nuvens na subida e o deserto na descida.}

O jeito mais direto de transportar calor: uma célula de convecção
gigante. Sobe no equador (ar quente e úmido \javimos{5}{convecção
profunda}), desce nos subtrópicos. Mas por que ela para em $30°$ e não
vai até o polo?

\begin{formadif}[momento angular fabrica o jato]
Anel de ar em repouso no equador tem momento angular $\Omega R^2$.
Levado a $\phi$ conservando
$M = (\Omega R\cos\phi + u)R\cos\phi$:
\[ u(\phi) = \Omega R\,\frac{\sin^2\phi}{\cos\phi}
\quad\Rightarrow\quad u(30°) = 134\ \mathrm{m/s}. \]
Observa-se $\sim$40\,m/s: turbilhões (os ciclones do Cap.~13!) roubam o
excesso. Ainda assim, o \textbf{jato subtropical} vive em $\sim$30°, a
12\,km — a célula de Hadley \emph{precisa} terminar onde o jato que ela
mesma fabrica se torna instável (T2).
\end{formadif}

\begin{noquadro}[o mapa-múndi climático em uma célula]
ramo que sobe (ZCIT): chuva diária equatorial \hfill
\javimos{7}{ITCZ no mapa de chuva}\\
ramo que desce ($30°$): Saara, Atacama, Kalahari — e as altas
semifixas do Atlântico/Pacífico Sul\\
ramo de retorno $+$ Coriolis \javimos{10}{Coriolis}: alísios de SE
(HS) / NE (HN)
\end{noquadro}

\section{Vento térmico: o jato polar}\label{sec:termico}

\quadro{Anunciar: ``a dedução mais elegante do curso — duas equações
que vocês já sabem, uma derivada, e sai o jato''. Fazer com calma.}

\begin{formadif}[relação do vento térmico]
Derivando o vento geostrófico \javimos{10}{geostrofia} na vertical e
usando hidrostática $+$ gás ideal \javimos{1}{as duas leis}:
\[ \frac{\partial u_g}{\partial z} = -\frac{g}{f\,T}\,
   \frac{\partial T}{\partial y}. \]
Leitura física via hipsométrica \javimos{1}{camada quente $=$ espessa}:
as superfícies de pressão se inclinam cada vez mais com a altura sobre a
zona de contraste — e inclinação de isóbara é vento.
\end{formadif}

\begin{noquadro}[a frase para decorar]
onde há contraste meridional de $T$, o vento de oeste \textbf{cresce
com a altura}\\
contraste concentrado (zona baroclínica) $\Rightarrow$ \textbf{jato
polar} na tropopausa, em cima da frente polar
\veremos{12}{frente polar}
\end{noquadro}

\begin{exemplo}[quanto jato dá o contraste?]
$\partial T/\partial y = -5$\,K por 1000\,km,
$f = -1{,}0\times10^{-4}$\,s$^{-1}$ ($45°$S), da superfície a 10\,km:
\[ \Delta u_g = \frac{g}{|f|\,T}
\left|\frac{\partial T}{\partial y}\right|\Delta z
 = \frac{9{,}8\times(5\times10^{-6})\times10^{4}}{10^{-4}\times260}
 \simeq 19\ \mathrm{m/s}. \]
Com contrastes de inverno e camadas mais profundas: os 40--80\,m/s
observados. \emph{Verificação no Caso~4 do notebook com perfil
realista.}
\end{exemplo}

Inverno: contraste maior $\Rightarrow$ jato mais forte e mais baixo em
latitude — por isso as tempestades de inverno
\veremos{13}{trilhas de tempestades}, e por isso a turbulência de céu
claro mora nas bordas do jato \javimos{5}{Richardson e CAT}.

\section{Ondas de Rossby}\label{sec:rossby}

\begin{noquadro}[vocabulário de vorticidade (Stull \S11.9)]
relativa $\zeta = \partial v/\partial x - \partial u/\partial y$
(o giro local) \quad planetária $f$ \quad absoluta $\zeta + f$\\
\textbf{potencial}: $(\zeta + f)/\Delta z$ — conservada seguindo o
escoamento adiabático:\\
coluna que estica gira mais; que achata, menos
\veremos{13}{esticamento e PV}
\end{noquadro}

\quadro{Desenhar cadeia de parcelas deslocada para o sul num planeta com
$\beta$: a vorticidade relativa muda para compensar $f$, o giro empurra
as vizinhas, a cadeia ondula e o padrão anda para oeste.}

\begin{formadif}[a restauração $\beta$]
Conservação de vorticidade absoluta $\zeta + f$ num deslocamento
meridional gera oscilação com dispersão (onda barotrópica):
\[ c = U - \frac{\beta}{k^2},\qquad
   \beta = \frac{df}{dy} = \frac{2\Omega\cos\phi}{R}. \]
A ``mola'' aqui não é o empuxo \javimos{5}{Brunt--Väisälä} — é a
variação de $f$ com a latitude: o planeta redondo é o meio elástico.
\end{formadif}

\begin{noquadro}[ler a dispersão]
ondas longas ($k$ pequeno): andam para \textbf{oeste} contra o jato\\
onda \emph{estacionária}: $\lambda_s = 2\pi\sqrt{U/\beta} \sim 6000$\,km
($U = 15$\,m/s)\\[2pt]
$\Rightarrow$ $\sim$4--5 meandros dando a volta no HS — confira em
qualquer carta de 50\,kPa \javimos{9}{cartas de altitude}
\end{noquadro}

As cristas e cavados dessas ondas são as altas e baixas móveis do dia a
dia \veremos{13}{ciclones e cavados}; quando a onda amplifica e
``quebra'', vêm os bloqueios e as ondas de calor/frio persistentes
\veremos{21}{extremos e bloqueios}. Dispersão e onda estacionária no
Caso~3 do notebook.

\section{O retrato completo}

\begin{noquadro}[os cinturões, de equador a polo]
\begin{tabular}{lll}
ZCIT ($\sim$0--10°) & convergência dos alísios & chuva diária \\
subtropical ($\sim$30°) & subsidência de Hadley & desertos, altas
semifixas \\
lat.\ médias (40--60°) & baroclinia $+$ Rossby & ciclones em fila \\
polar & subsidência fria & altas rasas \\
\end{tabular}
\end{noquadro}

Oceano: transporte de Ekman \javimos{10}{Ekman} $+$ rotacional do vento
médio $\Rightarrow$ giros subtropicais e correntes de contorno oeste
(Brasil, Golfo) — o outro caminhão de energia, com $\sim$1--2\,PW.

\textbf{Monções}: onde um continente tropical esquenta mais que o
oceano vizinho, forma-se uma baixa térmica sazonal que \emph{inverte}
os ventos entre verão e inverno — uma brisa marítima
\veremos{17}{brisas} em escala continental (Índia, oeste da África;
na América do Sul, o ``sistema de monção'' de verão que alimenta a
ZCAS). É a maior exceção regional ao retrato zonal dos cinturões.

\section{Síntese da aula}

\begin{noquadro}[mapa final --- canto do quadro]
\[ Q_{abs} - OLR \;\to\; T(\phi)\ (\sim5\,\mathrm{PW}) \;\to\;
   \text{Hadley},\ u = \Omega R\tfrac{\sin^2\phi}{\cos\phi} \]
\[ \frac{\partial u_g}{\partial z} =
-\frac{g}{fT}\frac{\partial T}{\partial y}\ \text{(jato)}
   \qquad c = U - \frac{\beta}{k^2}\ \text{(Rossby)} \]
\end{noquadro}

Daqui em diante o curso desce de escala: as massas de ar e frentes que
vivem nesses cinturões \veremos{12}{massas de ar e frentes}, os ciclones
que nascem da baroclinia \veremos{13}{ciclones extratropicais} e a
previsão numérica que integra tudo \veremos{20}{NWP}.

\section{Exercícios complementares}\label{sec:exercicios}

Soluções (professor) em \texttt{cap11\_solucoes.ipynb}.

\subsection*{Teóricos}

\begin{enumerate}[label=\textbf{T\arabic*.},itemsep=4pt]
  \item Escreva o balanço de energia por faixa de latitude e mostre que
        o transporte exigido é a integral do desequilíbrio; explique por
        que $T(90°) = 0$ é o teste de consistência do cálculo.
  \item Derive $u(\phi) = \Omega R\sin^2\phi/\cos\phi$ por conservação
        de momento angular e avalie em $10°, 20°, 30°$; discuta por que
        a célula de Hadley não pode ir além de $\sim$30°.
  \item Deduza a relação do vento térmico a partir de geostrofia $+$
        hidrostática ($+$ gás ideal) e mostre que no HS o jato fica a
        \emph{oeste} com ar frio ao sul.
  \item Pela conservação de $\zeta + f$, mostre que uma cadeia de
        parcelas deslocada meridionalmente oscila e propaga para oeste;
        obtenha $c = U - \beta/k^2$ (argumento heurístico basta).
  \item Estime o comprimento de onda estacionário para $U = 10, 15,
        25$\,m/s em $50°$S e conecte com o número de ondas planetárias
        observado ($\sim$4--5).
\end{enumerate}

\subsection*{Numéricos (no notebook)}

\begin{enumerate}[label=\textbf{N\arabic*.},itemsep=4pt]
  \item No Caso~1, adicione albedo alto ($A = 0{,}6$) além de $\pm70°$
        (gelo) e recalcule o transporte exigido; compare o pico.
  \item Com o EBM do \texttt{climlab} (Caso~2), varra a difusividade $D$
        e trace a posição da linha do gelo e o gradiente equador--polo
        contra $D$; interprete os limites $D\to0$ e $D\to\infty$.
  \item Trace $c(k)$ de Rossby para vários $U$ e ache numericamente o
        $k$ estacionário; verifique contra a fórmula do T5.
  \item Com o perfil zonal-médio de $T(y, z)$ fornecido no notebook,
        integre o vento térmico da superfície a 12\,km e localize o
        jato; compare inverno × verão.
\end{enumerate}

\vspace{1em}
\noindent\rule{\linewidth}{0.4pt}\\
{\scriptsize Material derivado de R.~Stull, \emph{Practical Meteorology}
(CC BY-NC-SA 4.0); este documento herda a mesma licença.}

\end{document}
